\section{Background}\label{s:background}

\subsection{Microservice Architecture Style}
In the microservice architecture style, applications are broken up into loosely-coupled independent components that communicate with each other and split the workload according to capability. Computation is focused in the services with communication between services being light-weight, often through APIs over HTTP. Due to this decentralized modular nature microservice systems are more scalable, robust and more easily modifiable \cite{mazzara_microservices_2018}. Since services are independent, bugs or faults are less likely to propagate. Furthermore, services can be independently scaled up and down and deployed, making it more flexible \cite{dragoni_microservices_2017}. The microservice architecture style does not necessarily decrease complexity, however. The strategies employed for managing services and the resources they run on can even lead to increased complexity and cost \cite{pirhonen_pains_2025}.

\subsection{Energy-efficient Microservices}
Studying and reducing the resource consumption, and particular the energy usage of the IT industry is a complex but crucial task \cite{verdecchia_green_2021}. The flexible and distributed nature of microservices form both opportunities and challenges in this regard. Scalability and flexibility mean that resources can be allocated resources according to demand, but this also requires in increase in complexity and overhead. Furthermore, techniques that may decrease energy usage, may reduce performance, reliability and availability. This is why new innovative techniques specifically designed for energy usage reduction in microservices are needed \cite{araujo_energy_2024, xiao_architectural_2024}.

\subsubsection{Measuring Energy Consumption of Microservice-based Systems}
In general, energy consumption measurement methods can be classified into three categories: measurement through external hardware, modelling power consumption based on resource utilization, and software solutions utilizing available hardware sensors to report energy consumption \cite{noureddine_review_2013}. In the microservices context, measurement of energy consumption can be divided into system-level and service-level measurement. As external hardware such as power meters only provide measures on a host- or system-level, they are not well-suited to microservice-based systems \cite{odea_measurement-driven_2025}. The method of modelling power consumption is employed (on a very basic level) by many simulators such as those discussed in this study, while software solutions are used to collect the ground truth these modelled values are compared to here. Examples of software energy measurement tools are Perf\footnote{\url{https://perfwiki.github.io/main}}, Scaphandre\footnote{\url{https://github.com/hubblo-org/scaphandre}}, Likwid\footnote{\url{https://github.com/RRZE-HPC/likwid}} and PowerJoular\footnote{\url{https://github.com/joular/powerjoular}} which all use the Running Average
Power Limit (RAPL) sensors which can report power consumption on the level of individual CPU sockets up to the entire system\cite{jay_experimental_2023}. Next to this, other software tools for measuring energy consumption of, for example, Python applications (e.g., EnergyScope\footnote{\url{https://www.energyscope.net/latest}}, PyJoules\footnote{\url{https://github.com/powerapi-ng/pyJoules}}) or GPUs (e.g., nividia-smi\footnote{\url{https://docs.nvidia.com/deploy/nvidia-smi/index.html}}) exists.

\subsection{Cloud Simulators}
Cloud simulators are used to simulate the behaviour of cloud datacenters and software or tasks running on them. Simulation can help with designing systems or applications, testing strategies and policies, and finding areas of improvement in current implementations. Pertinent topics of study at the moment are performance on large scales, security and resource usage in order to bring both environmental and monetary cost down \cite{mansouri_cloud_2020}. There is a substantial group of cloud simulators that support some form of energy usage simulation, but as these often originate out of scientific research, they are not always maintained after initial publication or cease being publicly available at all. Furthermore, not all cloud simulators are well suited to simulating microservices due to the containerized nature of the architecture style. Finally, the configurability of simulators can vary drastically. This can range from how fine-grained the user must define the system topology to how workloads can be defined in the simulator. \textit{CloudSim} \cite{buyya_energy_2024} is the most popular cloud simulator \cite{ahmed_cloud_2014}, a Java-based simulator that is relatively simple itself but for which extensions for \eg communication and GPU simulation are available and that other simulators such as \textit{CloudSim Plus} \cite{silva_filho_cloudsim_2017} are based on. \textit{GreenCloud} \cite{kliazovich_greencloud_2010} is a simulator first published in 2012 and one of the earliest simulators to include energy simulation capabilities for servers, but also network links and switches. It is not maintained any more. \autoref{tab:allSimulatorsAppendix} gives an overview of 67 simulators considered during the simulator selection phase of this experiment, which is described in \autoref{sec:simulators} along with a description of \textit{CloudSim}, \textit{DISSECT-CF}, \textit{NetworkCloudSim}, \textit{OpenDC} and \textit{SimGrid}, the simulators that were finally selected for this study.


